\documentclass[12pt]{article}
\def\activityid{F08-X-v2}\usepackage[letterpaper,margin=.65in,top=.60in,bottom=.65in]{geometry}
\usepackage[T1]{fontenc}
\usepackage[default]{sourcesanspro}
\usepackage{microtype,tikz,array,tabularx,fancyhdr,amsmath,amssymb}
\usepackage[hidelinks]{hyperref}
\usetikzlibrary{calc,arrows.meta}
\setlength{\parindent}{0pt}\setlength{\parskip}{7pt}
\setlength{\headheight}{0pt}\setlength{\footskip}{26pt}\setlength{\emergencystretch}{2em}
\pagestyle{fancy}\fancyhf{}\renewcommand{\headrulewidth}{0pt}\renewcommand{\footrulewidth}{.4pt}
\fancyfoot[L]{\scriptsize Bellingham Math Circle / Week 8 / \activityid}
\fancyfoot[R]{\scriptsize \thepage}
\definecolor{ink}{gray}{.12}\definecolor{soft}{gray}{.95}\color{ink}
\newcommand{\PageTitle}[2]{{\fontsize{10}{12}\selectfont\bfseries WEEK 8\quad STRATEGY LAB\hfill #1\par}\vspace{3pt}{\fontsize{26}{29}\selectfont\bfseries #2\par}\vspace{2pt}}
\newcommand{\NameLine}{{\small Name: \rule{2.65in}{.4pt}\hfill Date: \rule{1.35in}{.4pt}\par}\vspace{4pt}}
\newcommand{\Task}[2]{\par\vspace{3pt}{\large\bfseries #1}\par #2\par}
\newcommand{\Subhead}[1]{\par\vspace{3pt}{\large\bfseries #1}\par}
\newcommand{\RuleBox}[1]{\par\begingroup\setlength{\fboxsep}{8pt}\colorbox{soft}{\parbox{\dimexpr\linewidth-16pt\relax}{#1}}\endgroup\par}
\newcommand{\WriteBox}[2]{\par\textbf{#1}\par\vspace{2pt}\begin{tikzpicture}\draw[black!40,line width=.5pt] (0,0) rectangle (\dimexpr\linewidth-.5pt\relax,#2);\end{tikzpicture}\par}
\newcommand{\StudentFooter}{\vfill{\small Try it with objects. Explain by pointing, talking, or drawing.}\par}
\newcommand{\RookBoard}[3]{\begin{tikzpicture}[x=#3,y=#3]\draw[step=1,black!45] (0,0) grid (#1,#2);\draw[thick] (0,0) rectangle (#1,#2);\node[font=\Large] at ({#1-.5},.5){$\star$};\draw[-{Latex},thick] (0,-.35)--(#1,-.35);\draw[-{Latex},thick] ({#1+.35},#2)--({#1+.35},0);\end{tikzpicture}}
\newcommand{\Table}[3]{\begin{tikzpicture}[x=#3,y=#3]\draw[step=1,black!25] (0,0) grid (#1,#2);\draw[thick] (0,0) rectangle (#1,#2);\fill (0,0)circle(3pt);\draw[-{Latex},very thick] (0,0)--(.7,.7);\node[below] at ({#1/2},-.1){width #1};\node[rotate=90,above] at (-.2,{#2/2}){height #2};\end{tikzpicture}}
\newcommand{\GraphNodes}{\coordinate(A)at(0,0);\coordinate(B)at(3,0);\coordinate(C)at(1.5,2.6);\coordinate(D)at(1.5,.95);}
\newcommand{\Kfour}[1]{\begin{tikzpicture}[scale=#1]\GraphNodes\draw[very thick](A)--(B)--(C)--(A)--(D)--(B) (D)--(C);\foreach \n in {A,B,C,D}{\node[circle,draw,fill=white,minimum size=22pt,font=\bfseries]at(\n){\n};}\end{tikzpicture}}
\newcommand{\House}[1]{\begin{tikzpicture}[scale=#1]\coordinate(A)at(0,0);\coordinate(B)at(3,0);\coordinate(C)at(3,2);\coordinate(D)at(0,2);\coordinate(E)at(1.5,3.3);\draw[very thick](A)--(B)--(C)--(D)--(A) (D)--(E)--(C);\foreach \n in {A,B,C,D,E}{\node[circle,draw,fill=white,minimum size=22pt,font=\bfseries]at(\n){\n};}\end{tikzpicture}}

\begin{document}

\PageTitle{EXTRA / GRADES 6--7}{One new move, a new world}\NameLine
\RuleBox{Use TWO piles. You may remove any positive number from one pile, OR remove the SAME positive number from both. Take the last counter to win. This is Wythoff's game.}
\Task{1. Break the matching strategy.}{Why is \((4,4)\) no longer a trap? Classify all starts \((a,b)\) with \(0\leq a\leq b\leq7\). Work in order of \(a+b\); a trap has no move to an earlier trap. Treat \((0,0)\) as terminal.}
\Task{2. Build a smaller list.}{Start with \((0,0)\). Repeatedly choose the smallest nonnegative integer not yet in either column as \(a_n\), then put \(b_n=a_n+n\).}
\begin{center}\renewcommand{\arraystretch}{1.55}\begin{tabular}{c|*{6}{>{\centering\arraybackslash}p{.58in}}}n&0&1&2&3&4&5\\\hline$a_n$&0&1&&&&\\$b_n$&0&2&&&&\end{tabular}\end{center}
Does this list match your small-game traps? Explain why no rook move and no equal-take move can connect two listed pairs, even when pile order is swapped.
\WriteBox{What is proved? What still needs to be proved?}{.75in}
\Task{3. An unexpected number.}{Let \(\phi=(1+\sqrt5)/2\approx1.618\). A known theorem says the pairs are \((\lfloor n\phi\rfloor,\lfloor n\phi^2\rfloor)\), together with swapped pairs. Floors mean round down. Test \(n=1,\ldots,5\). Why does \(\phi^2=\phi+1\) explain the differences \(b_n-a_n=n\)?}
\textbf{Harder:} To prove the list is the whole losing set, what must you show about a position \emph{outside} it? Try finding such moves from \((2,4)\), \((5,8)\), and \((8,12)\).
\StudentFooter

\end{document}
